\uuid{DGyL}
\exo7id{4480}
\titre{exo7 4480}
\auteur{quercia}
\organisation{exo7}
\datecreate{2010-03-14}
\isIndication{false}
\isCorrection{false}
\chapitre{Série numérique}
\sousChapitre{Autre}
\cpge{mpsi-mp2i.analyse-asymptotique.a, pcsi.analyse-asymptotique.a, ptsi.analyse-asymptotique.a, mpsi-mp2i.procedes-sommatoires-discrets.a, pcsi.series-numeriques.a, ptsi.series-numeriques.a, mpsi-mp2i.procedes-sommatoires-discrets.b, pcsi.series-numeriques.b, ptsi.series-numeriques.b, bcpst1.limites-continuite.b, mp-mpi.series-numeriques-vectorielles.b, pc.series-numeriques, psi.series-numeriques.a, pt.series-numeriques.a, mp-mpi.integration-intervalle.e, bcpst2.series-reelles, bcpst2.integrales-generalisees}
\module{Analyse}
\niveau{L1}
\difficulte{}

\contenu{
\begin{enumerate}
    \item \question{Montrer qu'il existe $C\in\R$ tel que $\sum_{k=1}^n \frac{\ln k}k = \frac12\ln^2(n) + C +  o(1)$.}
    \item \question{Prouver~: $\frac{\ln2}2 -  \int_{t=1}^3\frac{\ln t}t\,d t \le C \le
               \frac{\ln2}2 + \frac{\ln3}3 -  \int_{t=1}^3\frac{\ln t}t\,d t$.}
    \item \question{Prouver~: $\sum_{k=1}^n \frac{\ln k}k = \frac12\ln^2(n) + C + \frac{\ln n}{2n} +  o\Bigl(\frac{\ln n}{n}\Bigr)$.}
\end{enumerate}
}
